% Part: first-order-logic
% Chapter: models-theories
% Section: expressing-props-of-structures

\documentclass[../../../include/open-logic-section]{subfiles}

\begin{document}

\olfileid{fol}{mat}{exs}

\olsection{\printtoken{P}{structure} चे गुणधर्म व्यक्त करणे}

\begin{explain}
फलने आणि संबंध यांच्यावरील अटी व्यक्त करणे, किंवा अधिक सर्वसाधारणपणे,
रचनेतील फलने आणि संबंध या अटींची पूर्ति करतात हे व्यक्त करणे अनेकदा
उपयुक्त आणि महत्त्वाचे असते. उदाहरणार्थ, एखाद्या भाषेसाठीच्या ज्या
!!{structure}s आपल्याला !!{predicate}s चा जो अर्थ ``पकडायचा'' आहे तो
पकडतात, त्या तसे न करणाऱ्या रचनांपासून वेगळ्या ओळखण्याचे मार्ग आपल्याला
हवे असतात. अर्थात, कोणती !!{structure}s आपल्याला ``अभिप्रेत'' आहेत हे
ठरवण्याचे पूर्ण स्वातंत्र्य आपल्याला आहे. उदाहरणार्थ, आपण असे ठरवू शकतो
की !!{predicate}~$\le$ चे अर्थनिर्धारण हा क्रमसंबंध असलाच पाहिजे; किंवा
आपल्याला फक्त~$\Lang L$ च्या अशा अर्थनिर्धारणांत रस आहे की ज्यांचा प्रांत
संचांचा बनलेला आहे आणि ज्यांत~$\Obj \in$ चे अर्थनिर्धारण ``चा
!!a{element} आहे'' या संबंधाने केलेले आहे. पण भाषेतील !!{sentence}s वापरून
आपण हे करू शकतो का? दुसऱ्या शब्दांत, !!a{structure}~$\Struct M$ वरील
कोणत्या अटी आपण~$\Struct M$ च्या भाषेतील !!a{sentence} (किंवा कदाचित
!!{sentence}s चा संच) वापरून व्यक्त करू शकतो? काही अटी आपल्याला व्यक्त करता
येणार नाहीत. उदाहरणार्थ, फक्त $\Domain M = \Nat$ असलेल्या
!!{structure}~$\Struct M$ मध्येच सत्य असे~$\Lang L_A$ चे कोणतेही वाक्य
नाही. ``प्रांतात फक्त नैसर्गिक संख्या आहेत'' हे आपण व्यक्त करू शकत नाही.
पण !!{structure}s चे काही ``रचनात्मक गुणधर्म'' आपण कदाचित व्यक्त करू शकतो.
!!{structure}s चे कोणते गुणधर्म आपण !!{sentence}s वापरून व्यक्त करू शकतो?
किंवा, दुसऱ्या प्रकारे विचारल्यास, एखादे !!a{sentence} (किंवा
!!{sentence}s चा संच) सत्य करणाऱ्या रचना म्हणून आपण
!!{structure}s च्या कोणत्या संग्रहांचे वर्णन करू शकतो?
\end{explain}

\begin{defn}[संचाचे प्रतिमान]
भाषा~$\Lang L$ मधील !!{sentence}s चा संच~$\Gamma$ असू द्या. प्रत्येक
$!A \in \Gamma$ साठी $\Sat{M}{!A}$ असेल, तर !!a{structure}~$\Struct M$
\emph{हे~$\Gamma$ चे प्रतिमान आहे} असे म्हणतो.
\end{defn}

\begin{ex}
वाक्य $\lforall[x][x \le x]$ हे~$\Struct M$ मध्ये तेव्हा आणि केवळ तेव्हाच
सत्य असते, जेव्हा $\Assign{\le}{M}$ हा परावर्ती संबंध असतो. वाक्य
$\lforall[x][\lforall[y][((x \le y \land y \le x) \lif x = y)]]$ हे
~$\Struct M$ मध्ये तेव्हा आणि केवळ तेव्हाच सत्य असते, जेव्हा
$\Assign{\le}{M}$ प्रतिसममित असतो. वाक्य
$\lforall[x][\lforall[y][\lforall[z][((x \le y \land y \le z)
      \lif x \le z)]]]$ हे~$\Struct M$ मध्ये तेव्हा आणि केवळ तेव्हाच सत्य
असते, जेव्हा $\Assign{\le}{M}$ संक्रमक असतो. म्हणून पुढील संचाची प्रतिमाने
\begin{align*}
\{\quad &\lforall[x][x \le x], \\
   & \lforall[x][\lforall[y][((x \le y \land y \le
    x) \lif x = y)]], \\
   &\lforall[x][\lforall[y][\lforall[z][((x \le y
      \land y \le z) \lif x \le z)]]] \quad \}
\end{align*}
म्हणजे नेमकी ती रचना होत ज्यांत~$\Assign{\le}{M}$ हा परावर्ती,
प्रतिसममित आणि संक्रमक, म्हणजेच अंशतः क्रम, असतो. त्यामुळे ही वाक्ये
आपण \emph{अंशतः क्रमांच्या प्रथम-क्रम उपपत्तीची} स्वयंसिद्धके म्हणून घेऊ
शकतो.
\end{ex}

\end{document}
